\section{The rational hitting point}
\label{sec:construction}

We now turn the integration operators into finite matrices and prove
Theorem~\ref{thm:main}. The hitting matrices depend only on the size parameters;
the representation of a polynomial is used to prove that they detect it.

\subsection{Taylor coefficients as a matrix space}

For $d\ge1$ and a characteristic-zero field $F$, let
\[
 \pi_d:F[[z]]\longrightarrow F^d,
 \qquad
 \pi_d\left(\sum_{r\ge0}a_rz^r\right)=(a_0,\ldots,a_{d-1})^{\mathsf T}.
\]
This identifies $F[[z]]/(z^d)$ with the column space having ordered basis
$1,z,\ldots,z^{d-1}$. The operator
$\mathcal J_i(a)=\int_0^z a(t)/(t+i)\,dt$ increases vanishing order, so it
induces a linear map on this quotient. Its matrix $T_i\in\Mat_d(\Q)$ is
\begin{equation}\label{eq:explicit-matrices}
 (T_i)_{r,q}=
 \begin{cases}
 \displaystyle\frac{(-1)^{r-q-1}}{r\,i^{\,r-q}},&0\le q<r<d,\\[4pt]
 0,&0\le r\le q<d.
 \end{cases}
\end{equation}
Indeed,
\[
 \frac1{z+i}=\sum_{a\ge0}\frac{(-1)^az^a}{i^{a+1}},
 \qquad
 \mathcal J_i(z^q)
 =\sum_{r>q}\frac{(-1)^{r-q-1}}{r\,i^{r-q}}z^r.
\]
Every denominator in the nonzero case is nonzero, since $r\ge1$ and $i\ge1$.
Discarding terms of degree at least $d$ before applying $\mathcal J_i$
does not change any retained coefficient. Thus, for every $a\in F[[z]]$,
\begin{equation}\label{eq:jet-intertwining}
 \pi_d(\mathcal J_i a)=T_i\pi_d(a).
\end{equation}

Let $e_0=\pi_d(1)$. For a word $w=i_1\cdots i_k$, repeated application of
\eqref{eq:jet-intertwining} gives
\[
 T_{i_1}\cdots T_{i_k}e_0=\pi_d(h_w).
\]
The rightmost matrix acts first, exactly as in
$h_w=\mathcal J_{i_1}\cdots\mathcal J_{i_k}(1)$.
For the empty word both products are identities. Consequently, for every
$f=\sum_w f_wx_w\in F\langle x_1,\ldots,x_n\rangle$,
\begin{equation}\label{eq:first-column}
 f(T_1,\ldots,T_n)e_0=\pi_d(h_f).
\end{equation}
This identity also accounts for constant terms, since a scalar is evaluated
as a scalar matrix.

\subsection{Detection and construction cost}

Suppose first that $ns\ge2$, and set $d=2ns^2$ in
\eqref{eq:explicit-matrices}. Let $F$ have characteristic zero, and let
$f\in\mathsf{Form}_{n,s}(F)$ be nonzero. By
Lemma~\ref{lem:word-independence}, its series $h_f$ is nonzero.
Proposition~\ref{prop:formula-system} represents this series as $uGv$
in a system of dimension $m\le2s$ satisfying the hypotheses of
Theorem~\ref{thm:multiplicity}. Hence
\begin{align*}
 \ord_0h_f
 &\le (n-1)\binom m2+m-1\\
 &\le (n-1)\binom{2s}{2}+2s-1
 \le 2ns^2-1=d-1.
\end{align*}
The last inequality follows from
\[
 (2ns^2-1)-\left((n-1)\binom{2s}{2}+2s-1\right)
 =s(2s+n-3)\ge0.
\]
Therefore $\pi_d(h_f)\ne0$. Equation~\eqref{eq:first-column} supplies a
nonzero first column of $f(T_1,\ldots,T_n)$, proving detection.

When $n=s=1$, a formula of size at most one is a single scalar or variable
leaf. The one-by-one tuple $T_1=[1]$ detects every nonzero polynomial in this
class. This gives dimension one in the remaining case.

For completeness, we give the bit cost of constructing the tuple.
With $d=2ns^2$, put
\[
 B=3+\lceil\log_2(d+1)\rceil
       +d\lceil\log_2(n+1)\rceil.
\]
Every nonzero entry in \eqref{eq:explicit-matrices} has numerator $1$ or $-1$
and positive denominator $r i^{r-q}$, of bit length at most $B$.
These fractions are already in lowest terms. For each $i$, precompute
$i,i^2,\ldots,i^{d-1}$ by repeated multiplication, then form each denominator
by multiplying one of these powers by $r$. There are $O(nd^2)$ such integer
operations on $O(B)$-bit integers. Schoolbook multiplication gives total
bit cost $O(nd^2B^2)$, with output length $O(nd^2B)$.
Both bounds are polynomial in $n,s$, and the stated entry bound in
Theorem~\ref{thm:main} follows by substituting $d=2ns^2$.
The branch $n=s=1$ has constant output and cost.
This completes the proof of Theorem~\ref{thm:main}.

\subsection{Acyclic path programs}

The formula-to-path conversion is the only step that uses the tree model.
The remaining argument applies to any strictly upper triangular coefficient
representation, with its number of states in place of the formula-size bound.

\begin{samepage}
\begin{proposition}
For every pair of integers $n,m\ge1$,\label{prop:path-programs}\label{rem:path-programs}
the matrices in
\eqref{eq:explicit-matrices}, with
\[
 d=(n-1)\binom m2+m,
\]
form a rational hitting point for the following class over every
characteristic-zero field $F$: nonzero polynomials
$f=\sum_w f_wx_w\in F\langle x_1,\ldots,x_n\rangle$ whose coefficients admit
a representation $f_w=uB_wv$, where $u\in F^{1\times m}$,
$v\in F^{m\times1}$, and $B_1,\ldots,B_n\in\Mat_m(F)$ are strictly upper
triangular. The tuple depends only on $n,m$ and can be constructed in time
polynomial in $n,m$, measured in bit operations.
\end{proposition}
\end{samepage}

\begin{proof}
For such a representation, define
$G=\sum_{|w|<m}B_wh_w$. Strict upper triangularity makes all longer products
zero. As in Proposition~\ref{prop:formula-system}, this gives
$G(0)=I_m$ and
\[
 G'=\left(\sum_{i=1}^n\frac{B_i}{z+i}\right)G,
 \qquad uGv=h_f.
\]
Thus $G$ is invertible over $F[[z]]$, and the system has the denominator
and numerator degrees in \eqref{eq:formula-system}.
The series $h_f$ is nonzero by Lemma~\ref{lem:word-independence}.
Theorem~\ref{thm:multiplicity} gives $\ord_0 h_f\le d-1$, so
\eqref{eq:first-column} proves detection. The preceding entry and output
cost calculation applies to this value of $d$ as well, and $d$ is polynomial
in $n,m$.
\end{proof}

In particular, an acyclic path program on $m$ vertices with scalar- or
variable-labelled edges has such a representation by the scalar-edge
elimination in Lemma~\ref{lem:formula-representation}. The same vertex bound
holds for the affine-linear edge labels allowed in an algebraic branching
program. To see this, order its vertices topologically and write each edge
label as $c_0+\sum_{i=1}^n c_i x_i$. Let $E_{ab}$ sum the coefficients $c_0$
over edges from $a$ to $b$, and let $(N_i)_{ab}$ sum their coefficients $c_i$.
Both $E$ and each $N_i$ are strictly upper triangular. Expanding the ordered
path products and putting $Q=(I_m-E)^{-1}$ gives, for source $\sigma$ and
sink $\rho$,
\[
 f_{i_1\cdots i_k}
 =e_\sigma^{\mathsf T}Q N_{i_1}Q\cdots N_{i_k}Qe_\rho.
\]
For the empty word the right side is $e_\sigma^{\mathsf T}Qe_\rho$.
Thus $u=e_\sigma^{\mathsf T}Q$, $B_i=N_iQ$, and $v=e_\rho$ give the
required representation without adding vertices.
